\uuid{mfpt}
\titre{Étude de suite}
\auteur{bodin ; hanani ; mzari}
\organisation{exo7}
\licence{CC BY-NC-SA 4.0}
\datecreate{2020-01-20}
\niveau{L1}
\module{Analyse}
\chapitre{Suite}
\sousChapitre{Autre}
\difficulte{3}
\theme{Suites, étude de suite}
\exo7qcmid{lille-1.suites.022}

\contenu{
\begin{question}
\qclassification{121.00}
\qtags{motcle=étude de suite}
Soit $\displaystyle u_n=\left|\frac{1}{n}-\frac{2}{n}+\frac{3}{n}-\dots+\frac{(-1)^{n-1}n}{n}\right|$. Quelles sont les bonnes réponses ?
\begin{answers}
    \bad{La suite $(u_n)$ est monotone.}
    \good{Les suites $(u_{2n})$ et $(u_{2n+1})$ convergent vers la même limite.}
    \bad{La suite $(u_n)$ est divergente.}
    \good{$\displaystyle \lim _{n\to +\infty}u_n=\frac{1}{2}$}
\end{answers}
\begin{explanations}
On vérifie que, pour tout $n\geq 1$,
$$u_{2n}=\frac{1}{2}\mbox{ et }u_{2n+1}=\frac{n+1}{2n+1}.$$
Donc les suites $(u_{2n})$ et $(u_{2n+1})$ convergent vers la même limite, à savoir $\displaystyle \frac{1}{2}$. D'après le théorème des suites extraites, la suite $(u_n)$ converge aussi vers $\displaystyle \frac{1}{2}$.
\end{explanations}
\end{question}
}
